%% kapa: exam=1 round=25 year=2014 period=2 form=A subject=appraisal_law subject_ko=부동산관계법규 seq=24 number=24
%% source: Q-Net 감정평가사 기출문제 (cst00309) artlSeq=5101207 file=extracted/1008/5101207/2014년 제25회 감정평가사 1차시험 기출문제/2교시(A형).hwp
%% conversion: kordoc 4.24.51 -> tools/gen1.py (auto); review=proofread-sonnet; flags=none
\begin{question}{24}
부동산 가격공시 및 감정평가에 관한 법령상 과징금의 부과 등에 관한 설명으로 옳은 것은?
\begin{choices}
  \item 감정평가업자가 위반행위로 취득한 이익의 규모를 고려하지 않고 과징금의 금액을 산정하여야 한다.
  \item 국토교통부장관은「부동산 가격공시 및 감정평가에 관한 법률」의 규정을 위반한 감정평가법인이 합병을 하는 경우 그 감정평가법인이 행한 위반행위는 합병 후 존속하거나 합병에 의하여 신설된 감정평가법인이 행한 행위로 보아 과징금을 부과ㆍ징수할 수 있다.
  \item 과징금납부의무자가 납부기한 내에 과징금을 납부하지 아니한 경우 국토교통부장관은 가산세를 징수할 수 있다.
  \item 과징금부과처분을 행정소송상 다투기 위해서는 소 제기에 앞서 반드시 행정심판을 거쳐야 한다.
  \item 과징금부과처분을 다투는 감정평가업자가 행정심판의 재결에 불복하는 경우 국토교통부장관에게 이의신청을 할 수 있다.
\end{choices}
\end{question}
\answer{2}
